\begin{tabular}{p{2.4cm}p{3.1cm}p{3.4cm}p{3.6cm}p{2.6cm}p{2.6cm}p{2.6cm}}
\toprule
Regime & Granularity & Stage taxonomy & Verification & Timeliness & Coverage & Public access \\
\midrule
\textbf{CEC energization tiers (California)} & MW by utility and tier, published as aggregates; project level only in SCE's public database (TN 266008, 268459) & three groups: signed agreement, active application, inquiry (2026 IEPR proposes a longer milestone list from inquiry to ramp) & utility-reported; no deposit or site-control test in the public record; confidence levels 70/33/0 or 100/50/10 applied by CEC staff & annual IEPR data request (December 2024, summer 2025, December 2025 vintages); SCE database twice a year & no size threshold; seven utilities reported for the 2025 IEPR (VEA's 2,600 MW are in Nevada) & docket documents and workshop decks; project-level detail confidential except SCE \\
\addlinespace
\textbf{ERCOT large load interconnection (Texas)} & MW by status, load zone, TSP, project type, size class and in-service year; monthly snapshots & no studies submitted; under ERCOT review; planning studies approved (Section 9.4/9.5 requirements met); approved to energize but not operational; observed energized & TSP planning studies reviewed by ERCOT; observed energization measured from telemetry (non-simultaneous peak); Batch Zero (July 2026) adds attestations and a load information form; PUCT rule 25.194 proposes USD 50,000/MW security & monthly (ERCOT Monthly, operational overview, NPRR1267 status report approved July 2025); board and legislative decks quarterly & loads of 75 MW or more at a point of interconnection (SB 6 threshold); 20 MW co-located with a resource & aggregated public reports; customer identities confidential \\
\addlinespace
\textbf{PJM load forecast adjustments} & MW by zone and year 2026-2046 (Tables B-9, B-9b); firm and non-firm by zone in the LAS summary & firm (electric service obligation or construction commitment) versus non-firm (letters of authorization, feasibility studies, other agreements) & utility submissions reviewed by PJM: 70\% utilization unless supported, ramps of at least 36 months, non-firm zero before 2030 and 50\% from 2030 with national scaling & annual cycle (requests July, review September to November, forecast January) & large load adjustments submitted by transmission owners; no MW threshold stated in the summary & load report, tables and the Load Analysis Subcommittee decks are public; project lists are not \\
\addlinespace
\textbf{EIA pilot data center surveys (federal)} & facility-level consumption for at least one data center per company (196 companies identified) & operating facilities only; measures energy use, energy sources, site characteristics, servers and cooling & voluntary web survey (Texas, Washington) and interviews (Northern Virginia and DC) & pilot launched March 25 2026; no publication schedule yet & three regions; California not included & results not yet published; EIA states the aim of faster cycles and finer detail \\
\addlinespace
\textbf{Texas SB 6 (2025)} & statute: disclosure duties per large load customer to the utility and ERCOT & interconnection standards to be set by the PUCT; PUCT rule 25.370 (March 2026) counts only loads with an executed interconnection agreement in the forecast after 2026 & disclosure of similar requests elsewhere in the state and of on-site backup generation; flat study fee of at least USD 100,000; proposed USD 50,000/MW financial security (rule 25.194) & standards and forecast rule in force from 2026; monthly ERCOT reporting continues & demand of 75 MW or more (the PUCT may lower it); curtailment service for loads of at least 75 MW during emergencies & utilities barred from selling or sharing customer submissions except to the PUCT and ERCOT \\
\addlinespace
\textbf{FERC Docket RM26-4 (federal)} & no data collection yet; the ANOPR (Oct 23 2025) proposes interconnection procedures for loads above 20 MW & study, agreement and transition provisions modelled on Order No. 2023 (CAISO comments); June 18 2026 show cause orders name five categories of reform & to be set in each RTO/ISO tariff; the orders give the six RTOs/ISOs 60 days to justify or file revisions (CAISO Docket EL26-71-000) & ANOPR Oct 2025; comments Nov 2025; FERC committed to act by June 2026 (Apr 16 2026) and issued the orders Jun 18 2026 & transmission-connected loads above 20 MW in the six RTO/ISO regions, about two-thirds of jurisdictional load & docket filings public through eLibrary (ferc.gov rejects non-browser clients, so only the CAISO and NERC filings and two news pages are in the raw store) \\
\addlinespace
\bottomrule
\end{tabular}